% GENERATED FILE. Edit canonical lessons, then run npm run generate:books.
% canonical-source-hash: fnv1a64:f4f88f92132074b1
% canonical-lessons: HI-C401-phatak, HI-C401-signal, HI-C401-janksan, HI-C401-laghna, HI-C401-takna

\chapter{Gates and Junctions}
\label{ch:hi-gates-and-junctions}
\hlchaptermodality{401}

\begin{chapteropening}
\textbf{By the end of this chapter:} \emph{I can talk about a level crossing, a signal and a junction, and gaze out at the fields.}

Complete the last lesson of chapter 401: I can talk about a level crossing, a signal and a junction, and gaze out at the fields.
\end{chapteropening}

\section[\texorpdfstring{\hi{फाटक} (phāṭak)}{फाटक (phāṭak)}]{\hi{फाटक} (phāṭak) — a gate; a level crossing}

\label{lesson:HI-C401-phatak}



\begin{quote}
Before the new one: say the Hindi for to check, then the Hindi for setting off.
\end{quote}

\subsection*{You'll want to know: \hi{फाटक}}
\textbf{\hi{फाटक}} — \emph{phāṭak} — \textquotedblleft{}a gate; a level crossing\textquotedblright{}.

\begin{quote}
\hi{फाटक} \hi{बंद} \hi{था}, \hi{इसलिए} \hi{हमें} \hi{रुकना} \hi{पड़ा।}
\end{quote}

\begin{quote}
\emph{phāṭak band thā, isliye hamẽ ruknā paṛā.}
\end{quote}

\begin{quote}
The crossing gate was shut, so we had to stop.
\end{quote}

Masculine.

\begin{cousinweb}
\textbf{\hi{फाटक}} (\emph{phāṭak}) goes back to \emph{*phāṭṭakka}, \textquotedblleft{}gate, door\textquotedblright{}, a form CDIAL reconstructs from its descendants (CDIAL 9069), which lists Hindi \emph{phāṭak}. Beside a railway line it is the barrier that drops across the road.
\end{cousinweb}

\begin{grammarlens}[title={What you've built}]
One more word for getting around.
\end{grammarlens}

\subsection*{Your turn}
\begin{itemize}
\raggedright
  \item \emph{Say it:} \emph{phāṭak}
  \item \emph{Say it:} \emph{phāṭak}, once more
  \item \emph{Say it:} \emph{phāṭak band thā, isliye hamẽ ruknā paṛā.} — The crossing gate was shut, so we had to stop.
  \item \emph{Recall:} say the Hindi for to check, then the Hindi for setting off, then say \emph{phāṭak} again
\end{itemize}

\begin{tcolorbox}[breakable,colback=teal!4,colframe=teal!35!black,title={Before you move on}]
What does \hi{फाटक} mean? (\textquotedblleft{}A gate; a level crossing\textquotedblright{}.) Say it once more.
\end{tcolorbox}
\section[\texorpdfstring{\hi{सिग्नल} (signal)}{सिग्नल (signal)}]{\hi{सिग्नल} (signal) — a signal}

\label{lesson:HI-C401-signal}



\begin{quote}
Before the new one: say the Hindi for setting off, then the Hindi for a gate.
\end{quote}

\subsection*{You'll want to know: \hi{सिग्नल}}
\textbf{\hi{सिग्नल}} — \emph{signal} — \textquotedblleft{}a signal\textquotedblright{}.

\begin{quote}
\hi{सिग्नल} \hi{लाल} \hi{था}, \hi{तो} \hi{ट्रेन} \hi{रुक} \hi{गई।}
\end{quote}

\begin{quote}
\emph{signal lāl thā, to ṭren ruk gaī.}
\end{quote}

\begin{quote}
The signal was red, so the train stopped.
\end{quote}

Masculine.

\begin{cousinweb}
\textbf{\hi{सिग्नल}} (\emph{signal}) is borrowed from English \emph{signal}, for the railway signal and the traffic light alike.
\end{cousinweb}

\begin{grammarlens}[title={What you've built}]
One more word for getting around.
\end{grammarlens}

\subsection*{Your turn}
\begin{itemize}
\raggedright
  \item \emph{Say it:} \emph{signal}
  \item \emph{Say it:} \emph{signal}, once more
  \item \emph{Say it:} \emph{signal lāl thā, to ṭren ruk gaī.} — The signal was red, so the train stopped.
  \item \emph{Recall:} say the Hindi for setting off, then the Hindi for a gate, then say \emph{signal} again
\end{itemize}

\begin{tcolorbox}[breakable,colback=teal!4,colframe=teal!35!black,title={Before you move on}]
What does \hi{सिग्नल} mean? (\textquotedblleft{}A signal\textquotedblright{}.) Say it once more.
\end{tcolorbox}
\section[\texorpdfstring{\hi{जंक्शन} (jaṅkśan)}{जंक्शन (jaṅkśan)}]{\hi{जंक्शन} (jaṅkśan) — a junction; a railway junction}

\label{lesson:HI-C401-janksan}



\begin{quote}
Before the new one: say the Hindi for a gate, then the Hindi for a signal.
\end{quote}

\subsection*{You'll want to know: \hi{जंक्शन}}
\textbf{\hi{जंक्शन}} — \emph{jaṅkśan} — \textquotedblleft{}a junction; a railway junction\textquotedblright{}.

\begin{quote}
\hi{यह} \hi{स्टेशन} \hi{एक} \hi{बड़ा} \hi{जंक्शन} \hi{है।}
\end{quote}

\begin{quote}
\emph{yah sṭeśan ek baṛā jaṅkśan hai.}
\end{quote}

\begin{quote}
This station is a big junction.
\end{quote}

Masculine.

\begin{cousinweb}
\textbf{\hi{जंक्शन}} (\emph{jaṅkśan}) is borrowed from English \emph{junction}, a station where railway lines meet; many station names end in it.
\end{cousinweb}

\begin{grammarlens}[title={What you've built}]
One more word for getting around.
\end{grammarlens}

\subsection*{Your turn}
\begin{itemize}
\raggedright
  \item \emph{Say it:} \emph{jaṅkśan}
  \item \emph{Say it:} \emph{jaṅkśan}, once more
  \item \emph{Say it:} \emph{yah sṭeśan ek baṛā jaṅkśan hai.} — This station is a big junction.
  \item \emph{Recall:} say the Hindi for a gate, then the Hindi for a signal, then say \emph{jaṅkśan} again
\end{itemize}

\begin{tcolorbox}[breakable,colback=teal!4,colframe=teal!35!black,title={Before you move on}]
What does \hi{जंक्शन} mean? (\textquotedblleft{}A junction; a railway junction\textquotedblright{}.) Say it once more.
\end{tcolorbox}
\section[\texorpdfstring{\hi{लाँघना} (lā̃ghnā)}{लाँघना (lā̃ghnā)}]{\hi{लाँघना} (lā̃ghnā) — to step over; to cross (a line)}

\label{lesson:HI-C401-laghna}



\begin{quote}
Before the new one: say the Hindi for a signal, then the Hindi for a junction.
\end{quote}

\subsection*{You'll want to know: \hi{लाँघना}}
\textbf{\hi{लाँघना}} — \emph{lā̃ghnā} — \textquotedblleft{}to step over; to cross (a line)\textquotedblright{}.

\begin{quote}
\hi{रेल} \hi{की} \hi{लाइन} \hi{लाँघना} \hi{ख़तरनाक} \hi{है।}
\end{quote}

\begin{quote}
\emph{rel kī lāin lā̃ghnā khatarnāk hai.}
\end{quote}

\begin{quote}
Crossing the railway line on foot is dangerous.
\end{quote}

Past, with \hi{ने} (he or she …): \hi{उसने} \hi{दीवार} \hi{लाँघी}, \emph{usne dīvār lā̃ghī}.

\begin{cousinweb}
\textbf{\hi{लाँघना}} (\emph{lā̃ghnā}) is a tadbhava of Sanskrit \emph{laṅghayati}, \textquotedblleft{}leaps over\textquotedblright{} (CDIAL 10905); the nasal before gh has become the nasal vowel ā̃.
\end{cousinweb}

\begin{grammarlens}[title={What you've built}]
One more word for getting around.
\end{grammarlens}

\subsection*{Your turn}
\begin{itemize}
\raggedright
  \item \emph{Say it:} \emph{lā̃ghnā}
  \item \emph{Say it:} \emph{lā̃ghnā}, once more
  \item \emph{Say it:} \emph{rel kī lāin lā̃ghnā khatarnāk hai.} — Crossing the railway line on foot is dangerous.
  \item \emph{Recall:} say the Hindi for a signal, then the Hindi for a junction, then say \emph{lā̃ghnā} again
\end{itemize}

\begin{tcolorbox}[breakable,colback=teal!4,colframe=teal!35!black,title={Before you move on}]
What does \hi{लाँघना} mean? (\textquotedblleft{}To step over; to cross (a line)\textquotedblright{}.) Say it once more.
\end{tcolorbox}
\section[\texorpdfstring{\hi{ताकना} (tāknā)}{ताकना (tāknā)}]{\hi{ताकना} (tāknā) — to look out (at); to stare}

\label{lesson:HI-C401-takna}



\begin{quote}
Before the new one: say the Hindi for a junction, then the Hindi for to step over.
\end{quote}

\subsection*{You'll want to know: \hi{ताकना}}
\textbf{\hi{ताकना}} — \emph{tāknā} — \textquotedblleft{}to look out (at); to stare\textquotedblright{}.

\begin{quote}
\hi{मीरा} \hi{खिड़की} \hi{से} \hi{खेतों} \hi{को} \hi{ताक} \hi{रही} \hi{थी।}
\end{quote}

\begin{quote}
\emph{Mīrā khiṛkī se khetoṁ ko tāk rahī thī.}
\end{quote}

\begin{quote}
Mira was gazing at the fields from the window.
\end{quote}

Past, with \hi{ने} (he or she …): \hi{उसने} \hi{आसमान} \hi{को} \hi{ताका}, \emph{usne āsmān ko tākā}.

\begin{cousinweb}
\textbf{\hi{ताकना}} (\emph{tāknā}) is a tadbhava of Sanskrit \emph{tarkayati}, \textquotedblleft{}guesses, thinks\textquotedblright{} (CDIAL 5716), which lists Hindi \emph{tāknā}, \textquotedblleft{}to look for\textquotedblright{}. How the sense moved to a long, fixed look is not spelled out there, so take the path as likely, not proven.
\end{cousinweb}

\begin{grammarlens}[title={What you've built}]
One more word for getting around.
\end{grammarlens}

\subsection*{Your turn}
\begin{itemize}
\raggedright
  \item \emph{Say it:} \emph{tāknā}
  \item \emph{Say it:} \emph{tāknā}, once more
  \item \emph{Say it:} \emph{Mīrā khiṛkī se khetoṁ ko tāk rahī thī.} — Mira was gazing at the fields from the window.
  \item \emph{Recall:} say the Hindi for a junction, then the Hindi for to step over, then say \emph{tāknā} again
\end{itemize}

\begin{tcolorbox}[breakable,colback=teal!4,colframe=teal!35!black,title={Before you move on}]
What does \hi{ताकना} mean? (\textquotedblleft{}To look out (at); to stare\textquotedblright{}.) Say it once more.
\end{tcolorbox}
